\section{Methods}
\label{sec:methods}

The analysis is split into independent workstreams, each writing a
\texttt{metrics.json} that feeds both its figure and the numbers quoted here, so a
single source produces every value.
The whole pipeline rebuilds from raw data with one driver script.

\subsection{Terrain}

We build a digital elevation model of the abbey promontory from AWS Terrain Tiles
(terrarium encoding, zoom 14, about \SI{6}{\meter} posting).
From it we compute the local slope and aspect within a \SI{400}{\meter} radius of
the tower and extract an elevation section along the dominant SW--NE axis.
The \SI{6}{\meter} posting under-resolves the sharp launch edge, so slope estimates
are reported as lower bounds pending airborne LIDAR.

\subsection{Wind climatology}

We take twenty years of ERA5 hourly \SI{10}{\meter} wind at the abbey (2004--2023,
\num{175320} hours) through the Open-Meteo archive
\autocite{hersbach2020era5,openmeteo_archive}.
A launch is counted favourable when the wind lies in the SW sector
(\ang{191}--\ang{259}, aligned with the escarpment aspect) at \SIrange{5}{12}{\meter\per\second}.
We report the favourable fraction annually and by meteorological season with
\SI{95}{\percent} confidence intervals taken across the twenty years.
Reconstructions of the medieval North Atlantic Oscillation differ on timing,
from a persistently positive Medieval Climate Anomaly mode
\autocite{trouet2009nao} to one with positive phases concentrated in later centuries
\autocite{ortega2015nao}, but neither points to the weak westerlies of a negative
phase, so modern reanalysis is a reasonable and probably conservative proxy
for Eilmer's epoch; the Twentieth Century Reanalysis provides an independent
stationarity check \autocite{slivinski2019twentieth}.

\subsection{Ridge lift}

We model orographic lift over the real terrain section with a first-order linear
solution in which the vertical air velocity follows the surface slope and decays
with height,
\begin{equation}
  w(x,z) = U\,\frac{\mathrm{d}h}{\mathrm{d}x}\,e^{-kz},
  \label{eq:ridge}
\end{equation}
where \(U\) is the incoming wind speed, \(\mathrm{d}h/\mathrm{d}x\) the along-wind
surface slope, \(z\) height above the surface, and \(k = 2\pi/L\) a decay rate set
by the escarpment horizontal scale \(L\).
We compare the resulting lift against a crude glider's sink rate to estimate the
depth of the soarable band.
\Cref{eq:ridge} is a deliberately conservative stand-in for a microscale large-eddy
or WAsP solve, which is deferred to a LIDAR-resolved follow-up.

\subsection{Glide}

Steady gliding flight sets the airspeed and sink rate.
For mass \(m\), wing area \(S\), trim lift coefficient \(C_L\), and air density
\(\rho\), the trim airspeed is \(V = \sqrt{2mg/(\rho S C_L)}\), the glide angle
\(\gamma\) satisfies \(\tan\gamma = 1/(\LD)\), and the still-air sink is
\(V\sin\gamma\).
With orographic upwash \(w\) the net sink is \(V\sin\gamma - w\), and the ground
range from launch height \(h\) is \(R = V\cos\gamma \cdot h/\max(V\sin\gamma-w,\,0)\),
extended without limit where the net sink is non-positive (the pilot soars).
We invert \(R = \SI{200}{\meter}\) to find the glide ratio required as a function of
wind speed, and compare it against the crude cambered-wing band (\LD\
\numrange{3}{7}) supported by low-Reynolds-number airfoil data
\autocite{lissaman1983lowre,hoerner1985lift,drela2014flight,faa_glider}.

\subsection{Pitch stability}

Longitudinal stability turns on the centre of gravity \(\cgfrac\) relative to the
neutral point \(\npfrac\), both in fractions of the mean chord.
A hanging pilot's CG lies well aft, between \num{0.35} and \num{0.55}\,\(c\).
A tailless cambered wing has its neutral point near the quarter chord
(\(\npfrac \approx 0.25\,c\)); a tail shifts it aft and adds pitch damping.
We linearise the short-period mode and form the state matrix in angle of attack and
pitch rate, whose largest real eigenvalue \(\lambda\) sets the growth rate; the time
for a pitch disturbance to double is \(t_2 = \ln 2/\lambda\).
We also integrate the nonlinear pitch response to an \ang{8} gust and record the
time to cross a \ang{60} loss-of-control threshold, comparing it against the flight
duration.
The treatment follows standard flight-dynamics theory \autocite{drela2014flight}.

\subsection{Monte Carlo synthesis}

The apparatus geometry and the launch day are unknown, so we sample the plausible
space with \num{40000} draws (fixed seed).
We draw wing area (\SIrange{5}{10}{\meter\squared}), all-up mass
(\SIrange{70}{90}{\kilogram}), trim \(C_L\) (\numrange{0.7}{1.1}), glide ratio
(\numrange{3}{7}), tower height (\SIrange{15}{25}{\meter}), escarpment slope
(\ang{8}--\ang{18}), ridge-lift efficiency (\numrange{0.5}{1.0}), CG
(\numrange{0.35}{0.55}\,\(c\)), aspect ratio (\numrange{3}{5}), and the neutral-point
shift from a tail (\numrange{0.18}{0.32}\,\(c\)).
Launch wind is drawn from the actual ERA5 SW-sector launchable speeds rather than
assumed, tying the synthesis to the measured climatology.
Each world is scored on whether the glide covers \SI{200}{\meter} and whether it
stays controlled (\(t_2\) longer than the flight) for the tailless and tailed
configurations; we report Wilson confidence intervals on the fractions.

\subsection{Impact biomechanics}

We treat ``broke both legs, survived, lame'' as data.
A feet-first impact decelerates the body over a stopping distance \(d\) made of leg
flexion and ground give, so the average total force is energy over distance,
\begin{equation}
  F = \frac{m v^2}{2 d}, \qquad \text{per leg } F/2,
  \label{eq:impact}
\end{equation}
for vertical impact speed \(v\).
The tibia fractures under axial and combined loading at
\SIrange{7.5}{11.3}{\kilo\newton} per leg, which sets a fracture-onset speed
\(v^\star = \sqrt{2 d F_\mathrm{tot}/m}\)
\autocite{salzar2014leg,abidin2026tibia,forman2012femur,ivarsson2005bending,femurFE2022}.
We compare a flared touchdown (vertical speed near the sink rate) against an
uncontrolled departure (vertical speed \(V\sin\theta\) for impact angle \(\theta\))
across the plausible range of ground give.
\Cref{eq:impact} is an average-force model, so the onset speeds are conservative
upper bounds on what the legs can take; real peak transients and bending loads would
lower the threshold.
